\section{A combinatorial lower bound for a fixed color}\label{sec:local}
Fix $z\in Q$ and put $S=Q\setminus\{z\}$. By
Lemma~\ref{lem:normalization}, we may work with a normalized strategy
without changing the matrices $M_i^z$. Throughout this section it is enough
to require winning only on the configurations in which $z$ occurs twice.
The argument proceeds from one fixed nonadjacent pair of $z$-vertices
to all five such pairs. We first obtain a pairwise inequality, then
analyze equality, and finally show that total weights 15 and 16 would
force incompatible equality cases.

\subsection{A missed-pair graph}
Fix two nonadjacent vertices $i,i+2$ and write a configuration in cyclic
order starting at $i$ as
\begin{equation}\label{eq:two-z}
(z,u,z,v,w),\qquad u,v,w\in S,\quad v\ne w.
\end{equation}
Set
\[
t=f_{i+1}(z,z),\qquad
\alpha(w)=f_{i+3}(z,w),\qquad
\beta(v)=f_{i+4}(v,z).
\]
Normalization gives $t\in S$, $\alpha(w)\in S\setminus\{w\}$, and
$\beta(v)\in S\setminus\{v\}$. Define
\begin{equation}\label{eq:missed}
E=\{(v,w)\in S^2:v\ne w,\ \alpha(w)\ne v,\ \beta(v)\ne w\}.
\end{equation}
The variables $u,v,w$ describe actual hat colors, whereas $t,\alpha,\beta$
are answers supplied by the strategy. The five observations give:
\begin{center}\small
\begin{tabular}{@{}cccl@{}}
\toprule
Player&Hat&Observation&Condition for a correct guess\\
\midrule
$i$&$z$&$(w,u)$&$M_i^z(w,u)=1$\\
$i+1$&$u$&$(z,z)$&$u=t$\\
$i+2$&$z$&$(u,v)$&$M_{i+2}^z(u,v)=1$\\
$i+3$&$v$&$(z,w)$&$\alpha(w)=v$\\
$i+4$&$w$&$(v,z)$&$\beta(v)=w$\\
\bottomrule
\end{tabular}
\end{center}
View $E$ as a bipartite graph whose row vertices are labeled by $v$ and whose
column vertices are labeled by $w$. Among the 12 ordered unequal pairs,
$\alpha$ guesses exactly four and $\beta$ guesses exactly four. Their sets
may intersect, so
\begin{equation}\label{eq:missed-size}
|E|\geq4.
\end{equation}
Every vertex of this graph has degree at most~3; in particular, one vertex
cannot cover all its edges. Notice that $E$ does not depend on~$u$.
Even when a row and a column have the same color label, they are distinct
vertices of this bipartite graph. A vertex cover is a set of row and
column vertices meeting every edge; it is not a set of colors.

For $u\ne t$, the middle player in~\eqref{eq:two-z} is wrong. For
$(v,w)\in E$, both end players are wrong as well. Thus one of the two
players wearing $z$ must be correct, which is equivalent to
\begin{equation}\label{eq:slice-cover}
M_i^z(w,u)=1\quad\text{or}\quad M_{i+2}^z(u,v)=1.
\end{equation}
For a fixed $u$, the selected columns $w$ and selected rows $v$ must
therefore form a vertex cover of~$E$.
For this fixed $u$, its size is exactly the number of selected entries
in column $u$ of $M_i^z$ plus the number in row $u$ of $M_{i+2}^z$.
The selected entries for different values of $u$ are disjoint. Thus the
three cover bounds may be added without double counting.

\begin{lemma}[Pair inequality]\label{lem:pair}
For every $i$,
\[
s_i(z)+s_{i+2}(z)\geq6.
\]
Consequently, $S(z)\geq15$.
\end{lemma}
\begin{proof}
For each of the three choices $u\in S\setminus\{t\}$, a vertex cover as
in~\eqref{eq:slice-cover} has at least two vertices by
\eqref{eq:missed-size}. Hence
\[
\sum_{w\in S}M_i^z(w,u)+\sum_{v\in S}M_{i+2}^z(u,v)\geq2.
\]
Sum over these three values of $u$ and include the nonnegative contribution
at $u=t$. This gives the pair inequality. Summing it over the five values
of $i$ counts every $s_i(z)$ twice and gives $2S(z)\geq30$.
\end{proof}

\subsection{All vertex covers of size two}
The equality case of Lemma~\ref{lem:pair} requires more information than
the number of edges in~$E$.

\begin{lemma}[Two-vertex covers]\label{lem:covers}
For a graph $E$ defined by~\eqref{eq:missed}, the following hold.
\begin{enumerate}
\item A cover consisting of a row $a$ and a column $b$ is possible only
when $a=b$. If it exists, it is the unique cover of size two, and $E$ is
the disjoint union of two copies of $K_{1,2}$ centered at that row and column.
\item There is at most one cover consisting of two rows and at most one
consisting of two columns.
\item If both pure covers exist, $E$ consists of all four edges of the
resulting $2\times2$ rectangle, and its row-label and column-label sets are
disjoint.
\end{enumerate}
\end{lemma}
\begin{proof}
Suppose that row $a$ and column $b$, with $a\ne b$, cover $E$. The submatrix
on $S\setminus\{a\}$ and $S\setminus\{b\}$ contains nine pairs. Exactly two
are diagonal pairs, since $|S\setminus\{a,b\}|=2$, leaving seven permitted
unequal pairs. None may belong to $E$. However, $\alpha$ can cover at most
one pair in each of the remaining three columns, and $\beta$ can cover at
most one in each of the three rows. Together they cover at most six, a
contradiction.

Now suppose that row $a$ and column $a$ cover $E$, and let
$T=S\setminus\{a\}$. The six unequal pairs in $T^2$ must all be covered
by $\alpha$ or $\beta$. Each function covers at most three such pairs.
Thus both cover exactly three, their covered sets are disjoint, and
$\alpha(T),\beta(T)\subseteq T$. In column $a$, the three legal rows lie
in $T$. The function $\beta$ guesses $a$ at none of these rows, while
$\alpha(a)\in T$ covers exactly one of the three pairs. Exactly two edges
of $E$ remain in column $a$. In row $a$, the function $\alpha$ never
answers $a$ in columns of $T$, and $\beta(a)\in T$ covers exactly one
of its three legal cells. Exactly two edges therefore remain in row $a$
as well. There are no other edges. The two stars have disjoint
vertex sets. A cover of a two-edge star that uses only one vertex must use
its center, so their two centers are the unique cover of size two.

If two distinct pairs of rows each cover $E$, every row incident with an
edge belongs to the intersection of the two pairs, which has size at most
one. This would give at most three edges, contradicting~\eqref{eq:missed-size}.
The same reasoning applies to columns. Finally, if a pair of rows and a
pair of columns both cover $E$, all edges lie in their $2\times2$ rectangle.
There are at least four edges, so every cell is an edge. A shared color
label between the two pairs would put a forbidden diagonal cell in this
rectangle. The label sets must therefore be disjoint.
\end{proof}

\subsection{Tight pairs and consecutive equalities}
The support of a Boolean matrix is called an $r\times s$ full rectangle
if it is the product of an $r$-element row set and an $s$-element column
set. In the following statement a rectangle with a zero dimension is empty.

\begin{lemma}[Structure of a tight pair]\label{lem:tight}
If $s_i(z)+s_{i+2}(z)=6$, then either
\begin{enumerate}
\item the supports of $M_i^z,M_{i+2}^z$ are full rectangles of sizes
$1\times3$ and $3\times1$, respectively; or
\item for some $k\in\{0,1,2,3\}$ they are full rectangles of sizes
$2\times k$ and $(3-k)\times2$, respectively.
\end{enumerate}
In particular,
\begin{equation}\label{eq:tight-counts}
(s_i(z),s_{i+2}(z))\in\{(0,6),(2,4),(3,3),(4,2),(6,0)\}.
\end{equation}
\end{lemma}
\begin{proof}
In the proof of Lemma~\ref{lem:pair}, the three slices $u\ne t$ must each
contribute exactly two entries, and the slice $u=t$ must contribute zero.
Each selected pair is therefore a cover of $E$ of size two.
There are no additional nonzero entries outside these covers: their
three contributions already account for the entire weight six.

If a mixed cover exists, Lemma~\ref{lem:covers} makes it unique: row $a$
and column $a$ for some $a\in S$. Thus $M_i^z(a,u)=1$ for all $u\ne t$,
and these are its only nonzero entries. Likewise, the only nonzero entries
of $M_{i+2}^z$ are $M_{i+2}^z(u,a)=1$ for $u\ne t$. This gives (i).

If there is no mixed cover, each slice chooses either the unique pure
column cover, when available, or the unique pure row cover. Let $U$ be the
set of slices using the column cover and put $k=|U|$. The nonzero entries
of $M_i^z$ have row labels equal to the two selected column labels of $E$
and column labels in $U$. The nonzero entries of $M_{i+2}^z$ have row
labels in $(S\setminus\{t\})\setminus U$ and column labels equal to the
two selected row labels of $E$. This proves (ii), including the cases in
which only one pure cover is available. When $k=0$ or $k=3$, the
unused cover need not exist; the corresponding matrix is empty, and
its zero-dimensional rectangle denotes just the empty support.
Counting entries gives
\eqref{eq:tight-counts}.
\end{proof}

\begin{corollary}[A restriction at two tight pairs]\label{cor:two-tight}
If $s_{i-2}(z)+s_i(z)=s_i(z)+s_{i+2}(z)=6$, then
$s_i(z)\in\{0,4\}$.
\end{corollary}
\begin{proof}
By~\eqref{eq:tight-counts}, the candidates are $0,2,3,4,6$. In the pair
$(i-2,i)$ the matrix $M_i^z$ is the right-hand matrix of
Lemma~\ref{lem:tight}, whereas in the pair $(i,i+2)$ it is the left-hand
matrix. For weights 2, 3, and 6, its required dimensions in these two roles
are, respectively,
\[
\begin{array}{c|ccc}
\text{weight}&2&3&6\\\hline
\text{right-hand role}&1\times2&3\times1&3\times2\\
\text{left-hand role}&2\times1&1\times3&2\times3
\end{array}
\]
A fixed nonempty matrix cannot have two different numbers of nonzero rows
or columns. These three weights are excluded, leaving only 0 and~4.
\end{proof}

\begin{proposition}[The lower bound 17]\label{prop:seventeen}
Every five-color winning strategy satisfies $S(z)\geq17$ for each $z\in Q$.
\end{proposition}
\begin{proof}
Put $a_j=s_{1+2j}(z)$ for $j\in\mathbb Z/5\mathbb Z$. Thus the player
order is $(1,3,5,2,4)$, and $a_j+a_{j+1}\geq6$. Lemma~\ref{lem:pair}
already gives $\sum_j a_j\geq15$.

If $\sum_j a_j=15$, all five pair inequalities are equalities. The equations
$a_{j+1}=6-a_j$ around an odd cycle force every $a_j=3$, contradicting
Corollary~\ref{cor:two-tight}.

If $\sum_j a_j=16$, put $e_j=a_j+a_{j+1}-6$. Summing gives
$\sum_j e_j=2\sum_j a_j-30=2$. Such a nonnegative integer vector has
either one entry 2 or two entries 1. In the latter case their positions
on a five-cycle are adjacent or nonadjacent; rotations and reflections
identify all choices of each kind. These are exactly three cases.
Alternating the five pair equations gives
\[
2a_0=(6+e_0)-(6+e_1)+(6+e_2)-(6+e_3)+(6+e_4),
\quad a_{j+1}=6+e_j-a_j.
\]
The following table keeps the same index choices in both columns:
\begin{center}\small
\begin{tabular}{@{}ccc@{}}
\toprule
$(e_0,e_1,e_2,e_3,e_4)$&$(a_0,a_1,a_2,a_3,a_4)$&Forbidden center\\
\midrule
$(2,0,0,0,0)$&$(4,4,2,4,2)$&$a_2=2$\\
$(1,1,0,0,0)$&$(3,4,3,3,3)$&$a_3=3$\\
$(1,0,1,0,0)$&$(4,3,3,4,2)$&$a_4=2$\\
\bottomrule
\end{tabular}
\end{center}
Each indicated center has both incident pair sums equal to six, but its
value is neither 0 nor 4. Corollary~\ref{cor:two-tight} excludes all
three cases. Thus 16 is impossible as well.
The count-preserving part of Lemma~\ref{lem:normalization} extends the
conclusion to strategies that were not initially normalized.
\end{proof}

\subsection{What the local bound does and does not prove}
The argument above uses only 240 of the 1020 proper colorings. The value
17 is sharp for that restricted requirement: Appendix~\ref{app:witness}
gives explicit tables winning on all two-occurrence colorings and having
$S(5)=17$. That same completion loses on other colorings. Thus a larger
lower bound for a full winning strategy must use additional information,
including compatibility with the zero- and one-occurrence families.
The next sections do this by keeping every proper coloring in the base
formula while fixing only the five counts. They do not freeze the
positions or the non-special answers of the local witness.

